\documentclass[11pt,a4paper]{article}

% =========================
% Language and encoding
% =========================
\usepackage[T2A]{fontenc}
\usepackage[utf8]{inputenc}
\usepackage[english,russian]{babel}

% =========================
% Page layout
% =========================
\usepackage[a4paper,margin=2.5cm]{geometry}
\usepackage{setspace}
\onehalfspacing

% =========================
% Typography and math
% =========================
\usepackage{lmodern}
\usepackage{microtype}
\usepackage{amsmath,amssymb,amsfonts}
\usepackage{mathtools}
\usepackage{bm}

% =========================
% Graphics and tables
% =========================
\usepackage{graphicx}
\usepackage{float}
\usepackage{booktabs}
\usepackage{array}
\usepackage{multirow}
\usepackage{tabularx}
\usepackage{longtable}
\usepackage{caption}
\usepackage{subcaption}

% =========================
% Links and references
% =========================
\usepackage[hidelinks]{hyperref}
\usepackage[nameinlink,noabbrev]{cleveref}

% =========================
% Algorithms and code
% =========================
\usepackage{algorithm}
\usepackage{algpseudocode}
\usepackage{listings}
\usepackage{xcolor}

% =========================
% Headers
% =========================
\usepackage{fancyhdr}
\pagestyle{fancy}
\fancyhf{}
\fancyhead[L]{GRA-Unified}
\fancyhead[R]{\thepage}

% =========================
% Bibliography
% =========================
\usepackage[numbers,sort&compress]{natbib}
\bibliographystyle{unsrtnat}

% =========================
% Custom commands
% =========================
\newcommand{\GRA}{\textsc{GRA}}
\newcommand{\Unified}{\textsc{Unified}}
\newcommand{\code}[1]{\texttt{#1}}
\newcommand{\vect}[1]{\bm{#1}}

% =========================
% Listings style
% =========================
\lstdefinestyle{mystyle}{
    basicstyle=\ttfamily\small,
    keywordstyle=\color{blue},
    commentstyle=\color{gray},
    stringstyle=\color{teal},
    numbers=left,
    numberstyle=\tiny,
    stepnumber=1,
    numbersep=8pt,
    showstringspaces=false,
    breaklines=true,
    frame=single
}
\lstset{style=mystyle}

% =========================
% Title
% =========================
\title{GRA-Unified: A Semantic Stability Framework for Modular AI Architectures}
\author{Your Name \and Coauthor Name}
\date{\today}

\begin{document}

\maketitle

\begin{abstract}
We present GRA-Unified, a modular semantic AI framework designed to support hierarchical reasoning, adaptive stability control, and composable agent coordination. The framework combines representation, control, and evaluation layers into a unified architecture for experimentation with advanced AI systems. We describe the design goals, internal modules, implementation strategy, and a preliminary evaluation protocol. The paper is intended as a technical foundation for further theoretical and experimental development.
\end{abstract}

\noindent\textbf{Keywords:} semantic AI, modular architecture, agent coordination, stability control, hierarchical reasoning, unified framework

\tableofcontents
\newpage

\section{Introduction}
The goal of GRA-Unified is to provide a common architecture for semantic AI experiments where reasoning, memory, control, and agent interaction are treated as distinct but interoperable layers. This paper documents the current design, terminology, and implementation skeleton of the framework. The structure is intentionally extensible so that future versions can add benchmark suites, stronger inference engines, and more explicit theory-to-code mappings.

This article is not a final claim of completeness. Instead, it is a reproducible technical description of a system under active development.

\section{Background}
Modern AI systems often combine multiple subsystems: embedding models, retrieval layers, planners, tools, and monitoring modules. In such systems, coupling between components can become difficult to reason about, especially when the goal is not only task performance but also semantic coherence and structural stability. GRA-Unified addresses this by organizing the system around explicit interfaces and higher-level control concepts.

\section{Framework Overview}
The framework is organized into the following conceptual layers:

\begin{itemize}
    \item \textbf{Semantic layer}: encodes meaning, symbols, and relations.
    \item \textbf{Control layer}: manages stability, gating, and state transitions.
    \item \textbf{Agent layer}: supports one or more autonomous reasoning units.
    \item \textbf{Memory layer}: stores persistent state, context, and traces.
    \item \textbf{Evaluation layer}: measures coherence, stability, and task utility.
\end{itemize}

Figure~\ref{fig:architecture} shows a simplified architecture diagram.

\begin{figure}[H]
    \centering
    \fbox{
    \begin{minipage}[c][5cm][c]{0.88\linewidth}
    \centering
    \vspace{1cm}
    \textbf{Insert architecture figure here}\\
    \vspace{0.5cm}
    Semantic Layer $\rightarrow$ Control Layer $\rightarrow$ Agent Layer $\rightarrow$ Evaluation Layer
    \vspace{1cm}
    \end{minipage}
    }
    \caption{Conceptual architecture of GRA-Unified.}
    \label{fig:architecture}
\end{figure}

\section{System Design}
The system is built around explicit state transitions and modular contracts. Each module exposes inputs, outputs, and internal parameters that can be traced independently. This makes it easier to validate stability properties and to compare different semantic policies under the same execution environment.

\subsection{Core Modules}
The core modules are:

\begin{enumerate}
    \item \textbf{Parser}: converts raw input into structured representations.
    \item \textbf{Semantic graph builder}: creates relational structures.
    \item \textbf{Stability controller}: adjusts internal thresholds and transitions.
    \item \textbf{Coordinator}: routes tasks among agents or submodules.
    \item \textbf{Logger}: records state, metrics, and event traces.
\end{enumerate}

\subsection{State Representation}
Let the global system state be denoted by \( S_t \). At time \( t \), the system can be written as
\[
S_t = \{m_t, c_t, a_t, r_t\},
\]
where \( m_t \) is memory state, \( c_t \) is control state, \( a_t \) is agent state, and \( r_t \) is the current reasoning trace.

\section{Implementation Sketch}
The implementation can be organized as a Python package with separate modules for semantics, control, agents, memory, and evaluation. A minimal repository layout could be:

\begin{lstlisting}[language=bash,caption={Example repository layout}]
GRA-Unified/
├── paper.tex
├── references.bib
├── src/
│   ├── gra/
│   │   ├── __init__.py
│   │   ├── semantic.py
│   │   ├── control.py
│   │   ├── agent.py
│   │   ├── memory.py
│   │   └── evaluation.py
├── figures/
│   ├── architecture.png
│   └── results.png
└── experiments/
    └── benchmark_runs.py
\end{lstlisting}

\section{Experimental Protocol}
To evaluate the framework, we recommend the following protocol:

\begin{itemize}
    \item Compare baseline and unified configurations on the same tasks.
    \item Measure coherence, response stability, and task completion quality.
    \item Report average values, variance, and failure modes.
    \item Include ablation tests for semantic and control modules.
\end{itemize}

Table~\ref{tab:metrics} provides an example template for reporting results.

\begin{table}[H]
    \centering
    \caption{Example evaluation metrics for GRA-Unified.}
    \label{tab:metrics}
    \begin{tabular}{lccc}
        \toprule
        Configuration & Coherence $\uparrow$ & Stability $\uparrow$ & Task score $\uparrow$ \\
        \midrule
        Baseline & 0.71 & 0.64 & 0.68 \\
        Unified w/o control & 0.75 & 0.59 & 0.70 \\
        GRA-Unified full & 0.83 & 0.81 & 0.79 \\
        \bottomrule
    \end{tabular}
\end{table}

\section{Results}
The results section should present empirical findings using consistent metrics and visualizations. If real benchmark data are available, insert plots such as convergence curves, ablation charts, or stability trajectories.

\begin{figure}[H]
    \centering
    \fbox{
    \begin{minipage}[c][5cm][c]{0.88\linewidth}
    \centering
    \vspace{1cm}
    \textbf{Insert results plot here}\\
    \vspace{0.5cm}
    Example: stability over time, benchmark comparison, or ablation chart
    \vspace{1cm}
    \end{minipage}
    }
    \caption{Example placeholder for a results figure.}
    \label{fig:results}
\end{figure}

\section{Discussion}
The main advantage of this framework is that it separates semantic organization from low-level execution. This improves interpretability and allows researchers to inspect where a failure originates. A second advantage is that the architecture is extensible: new reasoning modules can be added without rewriting the entire system.

A limitation is that the framework still requires concrete definitions of stability metrics and benchmarking tasks. Without these, the system remains a strong architectural proposal rather than a fully validated empirical platform.

\section{Conclusion}
GRA-Unified provides a structured basis for building semantic AI systems with explicit stability control and modular composition. The template in this paper is suitable for describing the repository, reporting experiments, and linking theory with implementation. Future work should include formal metrics, reproducible experiments, and richer comparisons with alternative architectures.

\section*{Acknowledgments}
Optional acknowledgments can be placed here.

\bibliography{references}

\end{document}